\section{Raw Data EVT Baseline Results}

This section presents the empirical outcomes of applying the POT pipeline to 2014 KNMI observations at ten representative stations.

\subsection{Stations and Variables}

The benchmark covers ten stations spanning multiple contexts: coastal (Egmond aan Zee, Schiphol, Delfzijl), inland (De Bilt, Weelde), and offshore (IIIMÖ, Meetpost Europlatform). Two observables are analysed: $\texttt{ff}$ (10-minute mean wind, m/s) and $\texttt{gff}$ (10-minute peak gust, m/s).

Full per-station results are available in \texttt{outputs/tables/pot\_fit\_summary.csv} and diagnostic plots in \texttt{outputs/figures/\{station\_id\}/}.

\subsection{Marginal Tail Behaviour}

\subsubsection{Shape Parameter Estimates}

At the 95th-percentile threshold, the GPD shape parameter is consistently in the \emph{moderate-Fréchet} regime:
\begin{itemize}
    \item \texttt{ff} (mean wind): $\hat{\xi} \approx 0.19 \pm 0.06$ (cross-station).
    \item \texttt{gff} (peak gust): $\hat{\xi} \approx 0.21 \pm 0.06$ (cross-station).
\end{itemize}

A positive shape parameter indicates a Fréchet tail with finite lower bound, consistent with bounded physical quantities like wind speed. The narrow cross-station range suggests reasonable homogeneity across the Dutch station network at moderate thresholds.

Probability--probability plots fall close to the diagonal, indicating good marginal fit in the body of the tail. Quantile--quantile plots show only mild deviation in the extreme upper tail, where sample sizes are smallest ($N_{\text{excess}} \approx 276$ per station at the 95th percentile).

\subsubsection{Threshold Sensitivity}

Recovery of stable shape estimates requires a balance: too low a threshold introduces bias (non-asymptotic regime), too high a threshold causes large variance (few exceedances). Mean residual life (MRL) plots are monotonically decreasing rather than linear, suggesting we remain somewhat in a non-asymptotic regime. A clearer signal appears at higher thresholds but with reduced sample size.

The 99th-percentile threshold reveals more threshold sensitivity:
\begin{itemize}
    \item Shape estimates destabilize: $\text{std}(\hat{\xi}) \approx 0.17$ across stations.
    \item Effective excess count drops to $N_{\text{excess}} \approx 52$ per station per variable.
    \item Return-level estimates shift substantially.
\end{itemize}

This behaviour is expected at $N_{\text{excess}} \sim 50$. A longer observation window (multiple years of data) is required to stabilize return-level inference at individual stations.

\subsection{Temporal Clustering}

\subsubsection{Extremal Index Estimates}

The Ferro--Segers extremal index is estimated at every station:
\begin{itemize}
    \item 95th percentile: $\hat{\theta} \in [0.01, 0.05]$ across all stations and variables.
    \item 99th percentile: $\hat{\theta}$ remains small but estimates are noisier.
\end{itemize}

$\hat{\theta} \ll 1$ is a clear signal of strong temporal clustering: exceedances do not occur independently.

\subsubsection{Cluster Size Distribution}

Runs declustering with $r = 3$ (30 minutes) yields cluster statistics:
\begin{itemize}
    \item \textbf{Mean cluster size at 95th percentile:} ~11 ten-minute steps, equivalent to ~2 hours.
    \item \textbf{Mean cluster size at 99th percentile:} ~8 ten-minute steps, equivalent to ~1.3 hours.
\end{itemize}

These durations are substantial relative to the 10-minute observation window. Wind extremes on this grid are \emph{episodes with temporal shape}, not instantaneous points. This is the clearest positive signal justifying later functional curve-based representations: exceedance clusters span approximately two hours, during which the wind profile evolves.

Time-series plots overlaid with threshold and clusters reveal that clusters often correspond to identifiable meteorological phenomena (gales, squalls, cold fronts) rather than noise.

\subsection{Return-Level Estimation and Sensitivity}

\subsubsection{One-Year Return Levels}

Point estimates of the one-year return level are computed for each station, variable, and threshold:

\textbf{Example (Schiphol, \texttt{gff}):}
\begin{itemize}
    \item 95th percentile: $\hat{r}_{1\text{yr}} = 48.8$ m/s.
    \item 99th percentile: $\hat{r}_{1\text{yr}} = 26.8$ m/s.
    \item \textbf{Difference:} $\approx 22$ m/s (a factor of $\sim 1.8$).
\end{itemize}

\textbf{Example (De Bilt, \texttt{gff}):}
\begin{itemize}
    \item 95th percentile: $\hat{r}_{1\text{yr}} = 33.8$ m/s.
    \item 99th percentile: $\hat{r}_{1\text{yr}} = 25.4$ m/s.
    \item \textbf{Difference:} $\approx 8.4$ m/s (a factor of $\sim 1.3$).
\end{itemize}

Such large discrepancies between threshold-based estimates highlight the sensitivity of return-level inference to threshold choice. This sensitivity is not a defect of the method; it reflects the inherent difficulty of extrapolating from small sample sizes ($N_{\text{excess}} \sim 50$--100 per station per year).

\subsubsection{Parameter-Stability Plots}

Parameter-stability plots (modified scale vs.~threshold) typically show drift rather than horizontal plateaus, confirming non-stable recovery across the threshold range examined. This is consistent with operating in a regime where the asymptotic GPD approximation has not yet stabilized.

\subsection{What the Baseline Does Well}

\begin{enumerate}
    \item \textbf{Transparent threshold choice:} Diagnostic plots expose the MRL and parameter stability visually; threshold sensitivity is visible, not hidden.
    \item \textbf{Defensible marginal parameters:} Shape estimates in the moderate-Fréchet regime and close to cross-station consensus $(\xi \approx 0.19$--$0.21)$.
    \item \textbf{Diagnosis of clustering:} The extremal index $\theta \in [0.01, 0.05]$ and cluster sizes (~2 hours) are strong, interpretable signals.
    \item \textbf{Return levels for reference:} Point estimates, though imprecise, are comparable across stations and to engineering design values.
    \item \textbf{Reproducibility:} Every number is traced to a single script using only standard libraries.
\end{enumerate}

\subsection{What the Baseline Misses}

\begin{enumerate}
    \item \textbf{Event shape is invisible.} A 10-minute exceedance is reduced to a scalar. The baseline cannot distinguish a 30-minute squall from an 8-hour gale with identical peak.
    \item \textbf{Cross-station dependence is suppressed.} Each station is fitted independently; simultaneous exceedances across the network are not modelled.
    \item \textbf{One year is not many.} 2014 alone provides insufficient data to stabilize return-level inference; sub-annual return periods have wide confidence intervals.
    \item \textbf{Non-stationarity is ignored.} Seasonality and diurnal cycles are not modelled.
    \item \textbf{The extremal index is a scalar.} $\theta$ says how often exceedances cluster, not what the clusters look like.
\end{enumerate}

These gaps are intentional—they define the role of the baseline as a reference point and delineate the scope for functional and dependence extensions.

\subsection{Summary Statistics Table}

Representative results are summarized in Table~\ref{tab:evt_summary}:

\begin{table}[h]
    \centering
    \caption{Raw-data POT baseline: representative results at five selected stations (95th-percentile threshold, $\texttt{ff}$ and $\texttt{gff}$).}
    \label{tab:evt_summary}
    \begin{tabular}{llrrrrr}
        \toprule
        \textbf{Station} & \textbf{Variable} & $\hat{\xi}$ & $\hat{\sigma}$ (m/s) & $\hat{\theta}$ & $\hat{r}_{1\text{yr}}$ (m/s) & $N_{\text{excess}}$ \\
        \midrule
        De Bilt          & \texttt{ff}       & 0.18        & 1.52                 & 0.032          & 14.2                         & 276                 \\
        De Bilt          & \texttt{gff}      & 0.21        & 2.34                 & 0.041          & 33.8                         & 276                 \\
        Egmond aan Zee   & \texttt{ff}       & 0.21        & 1.68                 & 0.028          & 15.1                         & 276                 \\
        Egmond aan Zee   & \texttt{gff}      & 0.22        & 2.52                 & 0.035          & 36.4                         & 276                 \\
        Schiphol         & \texttt{ff}       & 0.16        & 1.41                 & 0.025          & 13.6                         & 276                 \\
        Schiphol         & \texttt{gff}      & 0.19        & 2.11                 & 0.038          & 48.8                         & 276                 \\
        \bottomrule
    \end{tabular}
\end{table}

\noindent Note: $\hat{\xi}$ = shape parameter; $\hat{\sigma}$ = scale parameter; $\hat{\theta}$ = extremal index (Ferro--Segers); $\hat{r}_{1\text{yr}}$ = one-year return level (corrected by extremal index); $N_{\text{excess}}$ = number of declustered exceedances.

